% TOP_PROBLEM: {"id": 600005, "title": "Baker's Dozen — Completely periodic hyperbolic outer billiards", "category_id": 6, "category": {"id": 6, "name": "geometry", "display_name": "Geometry", "description": "Euclidean and non-Euclidean geometry, geometric structures.", "slug": "geometry", "order_index": 6, "created_at": "2026-07-31T15:26:25.671Z"}, "status": "open", "problem_number": "AMR-005-0005", "set_id": 13}
% FIRST_POSTED: 2026-10-01
% ENTRY_KIND: statement_only
% UnsolvedMath ID 600005; source catalogue code AMR-005-0005
% !TeX program = lualatex
% Definition and sources only; this file is not an LLM solution attempt.
\documentclass[11pt,a4paper]{article}
\usepackage[margin=25mm]{geometry}
\usepackage{fontspec}
\setmainfont{lmroman10-regular.otf}[BoldFont=lmroman10-bold.otf,
 ItalicFont=lmroman10-italic.otf,BoldItalicFont=lmroman10-bolditalic.otf]
\usepackage{amsmath,amssymb,mathtools}
\usepackage{unicode-math}
\setmathfont{latinmodern-math.otf}
\AtBeginDocument{\let\setminus\smallsetminus}
\usepackage{xurl,hyperref}
\hypersetup{colorlinks=true,urlcolor=blue}
\setcounter{secnumdepth}{0}
\setlength{\parindent}{0pt}
\setlength{\parskip}{5pt}
\setlength{\emergencystretch}{3em}
\begin{document}
\section*{Baker's Dozen — Completely periodic hyperbolic outer billiards}\label{600005}
{\small UnsolvedMath ID 600005 · AMR-005-0005 · Geometry · AMR Open Problem Lists}
\subsection{Definitions and mathematical statement}
For a compact convex geodesic polygon $P\subset\mathbb H^2$, orient the plane
and define the outer billiard map $T_P$ as follows. From a point $x$ outside
$P$, take the right supporting geodesic and its contact vertex $v$, and apply
the half-turn centered at $v$. The map is undefined when that support is
not a unique vertex. Write $\Omega_P$ for the set of exterior points whose
forward and backward iterates avoid all such singularities.

Call $P$ completely periodic when
\[
  \text{for every }x\in\Omega_P\text{ there is an integer }m(x)\ge1
  \text{ with }T_P^{m(x)}(x)=x.
\]
The requested task is a necessary and sufficient geometric characterization
of the polygons with this property, up to hyperbolic isometry. In particular,
the characterization should describe restrictions on side lengths, angles,
and their arrangement, rather than merely restating periodicity of the map.

The period may depend on the starting point; a common finite order for the
entire map is not part of the definition. Singular orbits are excluded, so
the classification does not depend on assigning arbitrary images at side
extensions. The question concerns all ordinary exterior orbits, not just
the dynamics induced on the circle at infinity.
\subsection{Short English statement}
Classify the hyperbolic polygons for which every nonsingular exterior
outer billiard trajectory eventually returns to its starting phase point.
\subsection{Sources}
\begin{itemize}
\item[S1] ULAM.AI and the UnsolvedMath Contributors, supplied problems.json, record 600005 (AMR-005-0005), AMR Open Problem Lists. Catalogue identity and source transcription; CC BY 4.0.
\url{https://huggingface.co/datasets/ulamai/UnsolvedMath}.
\item[S2] Tabachnikov - A Baker's Dozen of Problems (2015), Section 4 second unnumbered problem. Original source indexed by AMR-005-0005.
\url{https://amj.math.stonybrook.edu/html-articles/Files-2015-2024/14-01/}.
\end{itemize}
\end{document}
